% Hodina 46, 1. ročník SPŠ. Rozsah: Rozpis_vyuky_MAT_1_rocnik.tex, řádek 46.

% Plán: 1_rocnik/tematicke_plany/SPS.pdf; společná projektová šablona.

% Učebnice: ucebnice/01_Zakladni_poznatky/Algebraicke_vyrazy/Algebraické výrazy.tex

% Původ úloh: P2: učebnice, „Dělení mnohočlenu mnohočlenem“, přesný dělenec a dělitel; podmínka zde uvedena správně. Ostatní vlastní.

% Metodický zdroj: externi_zdroje/realisticky/Mnohocleny/05_Deleni_mnohoclenu.pdf

% URL: http://www.realisticky.cz/ucebnice/01%20Matematika%20S%C5%A0/01%20Z%C3%A1kladn%C3%AD%20poznatky/08%20Mnoho%C4%8Dleny/05%20D%C4%9Blen%C3%AD%20mnoho%C4%8Dlen%C5%AF.pdf

\documentclass[11pt,a4paper]{article}

% Spolecny styl pro ucitelske pripravy MAT (Work/MAT)
% Verze 1.1, 2026-10-23
\NeedsTeXFormat{LaTeX2e}
\ProvidesPackage{priprava}[2026/10/23 v1.1 Sablona ucitelskych priprav MAT]

% Zakladni typografie a cestina
\RequirePackage[T1]{fontenc}
\RequirePackage[utf8]{inputenc}
\RequirePackage[czech]{babel}
\RequirePackage{lmodern}
\RequirePackage{microtype}
\RequirePackage{amsmath,amssymb}
\RequirePackage{xcolor}
\RequirePackage{enumitem}
\RequirePackage[most]{tcolorbox}
\RequirePackage{tabularx,booktabs,array}
\RequirePackage{fancyhdr}
\RequirePackage{needspace}
\RequirePackage{geometry}
\RequirePackage[hidelinks]{hyperref}

\geometry{a4paper,top=17mm,bottom=18mm,left=23mm,right=23mm}
\definecolor{notegray}{gray}{0.39}
\setlength{\parindent}{0pt}
\setlength{\parskip}{2.5pt plus 1pt minus 0.5pt}
\setlist[itemize]{topsep=2pt,itemsep=1pt,parsep=0pt,leftmargin=1.45em}
\setlist[enumerate]{topsep=2pt,itemsep=2pt,parsep=0pt,leftmargin=1.65em}
\setcounter{secnumdepth}{0}
\renewcommand{\arraystretch}{1.12}

% Zahlavi: vlevo rocnik | globalni tema, vpravo cislo: nazev hodiny
% Zapati zustava zcela prazdne.
\pagestyle{fancy}
\fancyhf{}
\renewcommand{\headrulewidth}{0.35pt}
\renewcommand{\footrulewidth}{0pt}
\setlength{\headheight}{14pt}
\newcommand{\nastavPripravu}[4]{%
  \fancyhead[L]{\small #1\enspace\textbar\enspace #2}%
  \fancyhead[R]{\small #3: #4}%
}

% Minimalisticke cernobile boxy. Kazda definice zvlast.
\tcbset{commonlesson/.style={enhanced,breakable=false,colback=white,colframe=black,
  colbacktitle=white,coltitle=black,boxrule=0.45pt,
  arc=0.8mm,outer arc=0.8mm,left=3mm,right=3mm,top=1.6mm,bottom=1.6mm,
  toptitle=1mm,bottomtitle=0.7mm,fonttitle=\bfseries,
  attach boxed title to top left={xshift=2.5mm,yshift=-1.15mm},
  boxed title style={colback=white,colframe=white,boxrule=0pt,left=1mm,right=1mm}}}
\newtcolorbox{zadaniBox}[1][]{commonlesson,title={Zadání},#1}
\newtcolorbox{prikladBox}[1][]{commonlesson,title={Příklad},#1}
\newtcolorbox{definiceBox}[1][]{commonlesson,title={Definice},#1}

% Spolecny box pro zadani a jeho reseni, bez mezery mezi castmi.
% Pouziti:
%   \begin{prikladReseniBox}
%     Zadani ulohy.
%     \tcblower
%     \begin{align*} ... \end{align*}
%   \end{prikladReseniBox}
% Spodni cast ma jemny sedy podklad a delici linku; lze v ni sazet
% vice radku matematiky. Obe casti maji jen jeden spolecny ramecek.
\tcbset{solutionlesson/.style={
  commonlesson,
  enhanced jigsaw,
  breakable,
  colbacklower=black!3,
  segmentation style={solid,draw=black!25,line width=0.35pt},
  fontlower=\small\color{notegray},
  before lower={\noindent\textbf{Řešení:}\enspace},
  leftlower=3mm,rightlower=3mm
}}
\newtcolorbox{prikladReseniBox}[1][]{solutionlesson,title={Příklad},#1}
\newtcolorbox{zadaniReseniBox}[1][]{solutionlesson,title={Zadání},#1}

% Pomucky pro autora pripravy: sedive poznamky; \reseni pro starsi pripravy
\newcommand{\poznamka}[1]{\par\begingroup\small\color{notegray}\itshape
  \textnormal{\textbf{Poznámka pro učitele.}}\enspace #1\par\endgroup}
\newcommand{\reseni}[1]{\par\begingroup\small\color{notegray}
  \textbf{Očekávané řešení:}\enspace #1\par\endgroup}
\newcommand{\blok}[1]{\needspace{5\baselineskip}\subsection*{#1}}
\newcommand{\vyrok}[1]{\uv{#1}}


\nastavPripravu{1. ročník}{Algebraické výrazy}{46}{Dělení mnohočlenů: algoritmus}

\begin{document}

\blok{Výklad a souvislosti}

Dělení jednočlenem šlo po členech; mnohočlenný dělitel takto rozdělit nesmíme. Dlouhé dělení opakuje tři kroky: určit další člen podílu, vynásobit jím celý dělitel a odečíst vzniklý mnohočlen.

\begin{definiceBox}[title={Definice},breakable=false]
\textbf{Dělení mnohočlenů beze zbytku} znamená nalezení mnohočlenu $Q$ takového, že $A=BQ$, kde $B$ není nulový mnohočlen. Potom $A(x)/B(x)=Q(x)$ pro hodnoty s $B(x)\ne0$.
\end{definiceBox}


Dělenec i dělitel uspořádáme sestupně. Chybějící mocninu zapíšeme s nulovým koeficientem. Další člen podílu získáme dělením nejvyšších členů. Celý mezivýsledek odčítáme v závorce; jeho stupeň se tím sníží. Dnes vybíráme příklady s nulovým zbytkem.

\blok{Řešené příklady}

\begin{prikladReseniBox}[title={Příklad 1},breakable=false]
Vyděl $(x^2+3x+2):(x+1)$.
\tcblower
První člen podílu je $x$. Odečteme $(x^2+x)$ a zůstane $2x+2$. Další člen je $2$; po odečtení $(2x+2)$ zbývá $0$. Podíl je $x+2$ při $x\ne-1$. Kontrola: $(x+1)(x+2)=x^2+3x+2$.
\end{prikladReseniBox}


\begin{prikladReseniBox}[title={Příklad 2},breakable=false]
Vyděl $(2x^3-x^2-13x+5):(2x+5)$.
\tcblower
Postupné kroky:
\begin{align*}
(2x^3-x^2-13x+5)-x^2(2x+5)&=-6x^2-13x+5,\\
(-6x^2-13x+5)-(-3x)(2x+5)&=2x+5,\\
(2x+5)-1(2x+5)&=0.
\end{align*}
Podíl je $x^2-3x+1$. Pro hodnotu původního zlomku platí $x\ne-5/2$. Kontrola násobením vrací dělenec.
\end{prikladReseniBox}


\blok{Samostatné procvičení}

\begin{zadaniReseniBox}[title={Procvičení},breakable=false]
\begin{enumerate}[label=\alph*)]
\item Vyděl $(x^2-1):(x-1)$ dlouhým dělením.
\item Vyděl $(2x^2+5x+2):(2x+1)$.
\item Vyděl $(x^3+1):(x+1)$.
\end{enumerate}
\tcblower
\begin{enumerate}[label=\alph*)]
\item $x+1$ při $x\ne1$: nejprve odečteme $x(x-1)$, zůstane $x-1$, potom odečteme $1(x-1)$.
\item $x+2$ při $x\ne-1/2$. Po prvním kroku zbývá $4x+2$, po druhém $0$.
\item $x^2-x+1$ při $x\ne-1$. Do dělence doplníme nulové koeficienty u $x^2$ a $x$.
\end{enumerate}
\end{zadaniReseniBox}


\poznamka{Postup nejprve zapisuj celý, později lze použít obvyklé schéma. Učebnice u P2 chybně požaduje nenulový dělenec; správně se nesmí nulovat dělitel $2x+5$. Nenulový zbytek a jeho význam ponech další hodině.}

\end{document}
